\section{教学案例设计}

\subsection{案例背景与任务}

以某地级市（人口约300万）为研究区域，要求学生为该市规划一座"生态友好型城市物流园区"。物流园区选址是典型的多因子适宜性评价问题，涉及交通条件（高速公路距离、国省道距离）、地形地貌（坡度、地表起伏度）、生态环境（自然保护区缓冲区、重要水体距离、生态红线）、人口社会（居民点距离、耕地占用）及土地利用（地价水平）等因素的综合权衡。

案例的复杂性在于：各因子之间存在冲突（如交通便利区域往往地价较高），学生需通过科学的权重确定与叠加分析方法在多目标间寻求平衡。

\subsection{技术工具与数据}

\textbf{GIS平台}：ArcGIS Pro或QGIS作为空间分析主平台。\textbf{Python环境}：Jupyter Notebook结合ArcPy/GeoPandas实现空间数据处理。\textbf{生成式AI}：ChatGPT、文心一言或通义千问作为"智能协作者"，辅助代码生成与问题讨论。\textbf{数据源}：OpenStreetMap（道路网、居民点、水系）、地理空间数据云（DEM数据）、教师提供的区域基础地理数据。

\subsection{教学流程框架}

如图1所示，本案例采用"学生—AI—教师"三方互动的"人机协同四阶段"教学模式，各阶段层层递进、环环相扣。

\begin{figure}[htbp]
\centering
\fbox{图1 人机协同教学模式框架}
\caption{AI辅助GIS实践教学人机协同模式框架}
\label{fig:framework}
\end{figure}